\documentclass[11pt]{article}
\usepackage[T1]{fontenc}
\usepackage{amsmath,amssymb,geometry}
\geometry{margin=1in}
\title{Scalar tensor products do not add PI exponents\\Disproof of Conjecture 00000002751}
\author{ziangni-sys}
\date{October 8, 2026}
\begin{document}
\begin{center}
{\Large Scalar tensor products do not add PI exponents}\\[4pt]
Disproof of Conjecture 00000002751\\[4pt]
ziangni-sys --- October 8, 2026
\end{center}
\paragraph{Counterexample.}
Take the unital characteristic-zero PI algebras $A=B=\mathbb Q$ over
$\mathbb Q$. Their actual algebra tensor product is
$T=\mathbb Q\otimes_{\mathbb Q}\mathbb Q$. Multiplication defines an
algebra isomorphism
\[
 \phi:T\longrightarrow\mathbb Q,\qquad \phi(a\otimes b)=ab,
 \qquad \phi^{-1}(c)=1\otimes c.
\]
The inverse identity follows from the balancing relation
$a\otimes b=1\otimes ab$. In particular $T$ is a genuine nonzero
commutative algebra, and each of $A,B,T$ satisfies the nonzero
commutator polynomial $xy-yx$.

\paragraph{Codimensions from evaluations.}
Let $D$ be any nonzero unital commutative $\mathbb Q$-algebra, in
particular any of $A,B,T$. For $m\geq1$ let $P_m$ be the vector space
of formal noncommutative multilinear polynomials, with basis
$w_\sigma=x_{\sigma(1)}\cdots x_{\sigma(m)}$ for $\sigma\in S_m$.
Its actual evaluation map on all tuples in $D$ is
\[
 E_{D,m}:P_m\to\{D^m\to D\},\qquad
 E_{D,m}\Bigl(\sum_\sigma a_\sigma w_\sigma\Bigr)(t)
 =\sum_\sigma a_\sigma\prod_{j=1}^m t_{\sigma(j)}.
\]
The kernel is exactly $P_m\cap\operatorname{Id}(D)$. Commutativity
shows that
\[
 E_{D,m}(p)(t)=\Bigl(\sum_\sigma a_\sigma\Bigr)g_m(t),
 \qquad g_m(t)=\prod_{j=1}^m t_j.
\]
The function $g_m$ is nonzero since its value at $(1,\ldots,1)$ is
$1_D\ne0$. It belongs to the image because it evaluates any basis
word. Hence the image is precisely its one-dimensional $\mathbb Q$
span. The first isomorphism theorem gives the genuine codimensions
\[
 c_m(D)=\dim_{\mathbb Q}\bigl(P_m/(P_m\cap\operatorname{Id}(D))\bigr)
       =\dim_{\mathbb Q}\operatorname{im}E_{D,m}=1.
\]
Consequently the defining nth-root limits exist and
\[
 \operatorname{exp}(A)=\operatorname{exp}(B)=\operatorname{exp}(T)
 =\lim_{m\to\infty}\sqrt[m]{1}=1.
\]
Thus the conjectured exact sum would assert $1=1+1=2$, a
contradiction. The extra mixed-characteristic correction does not
repair a false assertion already in characteristic zero. We only
need this false clause to disprove the conjecture.

\paragraph{Formal certificate.}
Lean constructs the actual algebra tensor product and its standard
multiplication algebra equivalence. The multilinear evaluation map,
its quotient-by-kernel linear equivalence, one-dimensional image,
codimensions and root limit are established for a general nonzero
commutative $\mathbb Q$-algebra and instantiated separately on the
scalar factors and the tensor algebra. No exponent or codimension
sequence is stipulated. Final theorem audits use only standard Lean
and Mathlib axioms.
\end{document}
